\subsection{Non-Interactive Key Exchange (NIKE)}
\label{enike}

Our generic scheme additionally builds upon a Non-Interactive Key Exchange
(\ENIKE)~\cite{Bernstein06,CashKS08,FreireHKP13} scheme.
We additionally define a verification algorithm that checks whether a tuple is in fact a valid keypair.


\begin{definition} \label{defn:enike}
  A \defn{Non-Interactive Key Exchange (\ENIKE)} $\enike$ with verification is the tuple
  $\enike = (\keygen, \verifykeypair, \getsharedkey)$, where

\begin{itemize}
  \item $\keygen(\lambda) \rightarrow (\sk, \pk)$: Generate a secret key $\sk$ and corresponding public key $\pk$ with respect to the security parameter $\lambda$.
  \item $\verifykeypair(\sk, \pk) \rightarrow \zo$: Verify that $(\sk,\PK)$ is a valid output from $\enike.\keygen$.
  \item $\getsharedkey(\sk_1, \pk_2) \rightarrow \sharedkey$: Output the shared key $\sharedkey$ generated by the combination of one party's secret key $\sk_1$ and the other party's public key $\pk_2$.
  \item For correctness, we require that, for all $(\SK_1, \PK_1) \rgets \enike.\keygen(\lambda), (\SK_2, \PK_2) \rgets \enike.\keygen(\lambda)$, the following conditions hold:
\begin{align*}
  & \enike.\verifykeypair(\sk_i, \pk_i) = 1, \text{ for } i \in \{1,2 \} \text{ and} \\[1mm]
  & \enike.\getsharedkey(\sk_1, \pk_2) = \enike.\getsharedkey(\sk_2, \pk_1)
\end{align*}
\end{itemize}
\end{definition}

In addition to the requirement that a \ENIKE be \emph{session-key unrecoverable},
we additionally require that a \ENIKE be \emph{binding}.
We discuss both properties next.

\paragraph{Session-Key Recovery.}
A  \ENIKE that is \emph{session-key unrecoverable} ensures that given the public keys of two parties $\PK_1, \PK_2$,
it should be hard for an adversary to obtain the shared secret.
While prior security notions for non-interactive key exchange exist in the literature~\cite{CashKS08,FreireHKP13},
these notions allow an adversary to act as an active participant,
and assume a distinguishing game.
Our constructions require only the more restricted model where the adversary learns key material but does not contribute itself,
and must compute the shared secret directly.
We present this restricted notion in Figure~\ref{fg:enike-recovery}.

\begin{figure}[t]
  \begin{pchstack}[boxed,center,space=1em]
  \procedure[linenumbering]{$\UnrecoverabilityExp_{\enike,\adv}(\lambda)$}{
    (\sk_1, \pk_1) \rgets \enike.\keygen(\lambda) \\
    (\sk_2, \pk_2) \rgets \enike.\keygen(\lambda) \\
    \sharedkey \gets \enike.\getsharedkey(\sk_1, \pk_2) \\
    \sharedkey' \rgets \adv(\pk_1, \pk_2) \\
    \pcreturn 1\ \pcif \sharedkey' = \sharedkey \\
    \pcreturn 0
  }
\end{pchstack}

  \caption{Unrecoverability experiment defining the advantage of an adversary $\adv$ against a non-interactive key exchange (\ENIKE) $\enike$.
  }
\label{fg:enike-recovery}
\end{figure}



The advantage of an adversary $\adv$ against $\enike$ in the unrecoverability experiment as defined in Figure~\ref{fg:enike-recovery} is
\[
  \advant{rec}{\enike,\adv}(\lambda) = \Pr[\UnrecoverabilityExp_{\enike,\adv}(\lambda) = 1 ]
\]

\begin{definition}
A \ENIKE is \defn{session-key unrecoverable} if for all probabilistic polynomial time adversaries $\adv$,
  the function $\advant{rec}{\enike, \adv}(\lambda)$  is negligible.
\end{definition}


\paragraph{Binding.}
We introduce an additional security property that must hold when the adversary is allowed to generate both sides of the key exchange.
Informally, a \ENIKE is \emph{binding} when $\enike.\getsharedkey(\sk_1, \pk_2)$ must equal $ \enike.\getsharedkey(\sk_2, \pk_1)$,
even when the adversary is allowed to generate $(\sk_1, \pk_1), (\sk_2, \pk_2)$ itself.
The only restriction to the adversary is that both keypairs must be valid.
We formalize this requirement in Figure~\ref{fg:enike-binding}.

\begin{figure}[t]
  \begin{pchstack}[boxed,center,space=1em]
    \procedure[linenumbering]{$\BindingExp_{\enike,\adv}(\lambda)$}{
    (\sk_1^*, \pk_1^*, \sk_2^*, \pk_2^*) \rgets \adv(\lambda) \\
    \pcreturn 0\ \pcif \enike.\verifykeypair(\sk_1^*, \pk_1^*) \neq 1 \\
    \pcreturn 0\ \pcif \enike.\verifykeypair(\sk_2^*, \pk_2^*) \neq 1 \\
    %
    \pcreturn 1\ \pcif \enike.\getsharedkey(\sk_1^*, \pk_2^*) \neq \enike.\getsharedkey(\sk_2^*, \pk_1^*) \\
    \pcreturn 0
  }
  \end{pchstack}

  \caption{Binding experiment defining the advantage of an adversary $\adv$ against a non-interactive key exchange (\ENIKE) $\enike$.
  Here, the adversary can perform all $\enike$ operations directly.
  }
\label{fg:enike-binding}
\end{figure}

The advantage of an adversary $\adv$ against $\enike$ in the binding experiment as defined in Figure~\ref{fg:enike-binding} is
\[
  \advant{bind}{\enike,\adv}(\lambda) = \Pr[\BindingExp_{\enike,\adv}(\lambda) = 1 ]
\]

\begin{definition}
A \ENIKE is \defn{binding} if for all probabilistic polynomial time adversaries $\adv$,
the function $\advant{bind}{\enike, \adv}(\lambda)$  is negligible.
\end{definition}



\subsubsection{A Concrete \ENIKE.} \label{enike:concrete}
We now describe a concrete \ENIKE which is a non-interactive Diffie--Hellman key exchange,
with additionally a verification check to ensure that $\pk = g^{\sk}$.

\begin{itemize}
  \item $\enike.\keygen(\lambda) \rightarrow (\sk, \pk)$: Sample secret key $\sk \rgets \Zq$; derive public key $\pk\in \GG$ as $\pk = g^\sk$.
	    Output $(\sk, \pk)$.
  \item $\enike.\verifykeypair(\SK, \PK) \rightarrow \zo$: Output $1$ if $\PK = g^\SK$; otherwise, output $0$.
  \item $\enike.\getsharedkey(\sk_1, \pk_2) \rightarrow \sharedkey $: Derive the shared key $\sharedkey \in \GG$ as $\sharedkey \gets \pk_2^{\sk_1}$.
    Output $\sharedkey$.
\end{itemize}

We demonstrate in Appendix~\ref{appendix:enike-proofs} that
the concrete scheme is unrecoverable assuming the Computational Diffie-Hellman (CDH) problem is hard,
and is unconditionally binding.

